\documentclass[11pt]{article}

\usepackage[margin=1in]{geometry}
\usepackage{amsmath,amssymb,amsthm,mathtools}
\usepackage{enumitem}
\usepackage{hyperref}

\newtheorem{theorem}{Theorem}
\newtheorem{proposition}[theorem]{Proposition}
\newtheorem{lemma}[theorem]{Lemma}
\newtheorem{corollary}[theorem]{Corollary}
\theoremstyle{definition}
\newtheorem{definition}[theorem]{Definition}
\newtheorem{assumption}[theorem]{Assumption}
\newtheorem{remark}[theorem]{Remark}

\newcommand{\PM}{\mathrm{PM}}
\newcommand{\Pa}{\mathrm{Pa}}
\newcommand{\R}{\mathbb{R}}
\newcommand{\E}{\mathbb{E}}
\newcommand{\Prob}{\mathbb{P}}
\newcommand{\eps}{\varepsilon}

\title{Scaling Invariance of the Angular Lower-Tail Coefficient \\
$\kappa_L(\delta,\gamma;C)$, with Verification for the \\
Pareto-Mixture Copula $C_{\PM}$}
\author{}
\date{}

\begin{document}
\maketitle

\section{Setting and Statement of the Result}

Let $C:[0,1]^2\to[0,1]$ be a bivariate copula. For parameters
$\gamma>0$ and $\delta\in(0,2)$, define
\begin{equation}\label{eq:kappa-def}
\kappa_L(\delta,\gamma;C)
\;:=\;
\lim_{u\downarrow 0}
\frac{\log C\!\bigl(\gamma u^{\delta},\,\gamma^{-1}u^{2-\delta}\bigr)}{\log u},
\end{equation}
whenever the limit on the right-hand side exists in $\R$. This is an
angular generalization of the Ledford--Tawn lower-tail coefficient:
the curve $u\mapsto(\gamma u^{\delta},\gamma^{-1}u^{2-\delta})$
approaches $(0,0)$ along a direction in the unit square parametrized
by $\delta\in(0,2)$, while $\gamma>0$ rescales the two coordinates
inversely.

The intuition motivating the definition is that, when the limit
exists, the tail behavior captured by $\kappa_L$ should depend on the
\emph{direction} $\delta$ but not on the \emph{scale} $\gamma$. The
goal of this note is to formalize that intuition.

For notational convenience, throughout the paper we write
\begin{equation}\label{eq:Edef}
E_\gamma(u;d)
\;:=\;
\frac{\log C\!\bigl(\gamma u^{d},\,\gamma^{-1}u^{2-d}\bigr)}{\log u},
\qquad u\in(0,1),\;d\in(0,2),
\end{equation}
so that $\kappa_L(d,\gamma;C)=\lim_{u\downarrow 0}E_\gamma(u;d)$
whenever this limit exists.

\subsection{Assumptions}

Fix a copula $C$ and a reference scale $\tilde\gamma>0$. We state the
regularity conditions used in the main theorem up front, as the
analysis requires.

\begin{assumption}\label{ass:limit}
For every $d\in(0,2)$ the limit
$\tilde\kappa(d):=\lim_{u\downarrow 0}E_{\tilde\gamma}(u;d)$ exists in
$\R$. In particular, fixing the direction $\delta\in(0,2)$ of interest,
$\tilde\kappa(\delta)$ is finite.
\end{assumption}

\begin{assumption}[Local uniform convergence in $d$]\label{ass:uniform}
For every $\delta_0\in(0,2)$ and every $\eps>0$ there exist
$r=r(\delta_0,\eps)>0$ and $u_0=u_0(\delta_0,\eps)\in(0,1)$ such that
\[
\bigl|E_{\tilde\gamma}(u;d)-\tilde\kappa(d)\bigr|\;\le\;\eps
\qquad
\text{for all }u\in(0,u_0]\text{ and all }d\in(0,2)\text{ with }
|d-\delta_0|\le r.
\]
\end{assumption}

\begin{assumption}[Continuity of $\tilde\kappa$ at $\delta$]\label{ass:cont}
The function $d\mapsto\tilde\kappa(d)$ is continuous at $\delta$.
\end{assumption}

Assumptions~\ref{ass:limit}--\ref{ass:cont} are standard regularity
conditions in lower-tail dependence analysis. They hold for the
familiar copula models used in practice, including all examples of
the angular Ledford--Tawn framework whose lower tail is regularly
varying in $(u_1,u_2)$. They are exactly what we need for the
$\eps$-argument that follows.

\subsection{Main result}

\begin{theorem}\label{thm:main}
Under Assumptions~\textnormal{\ref{ass:limit}--\ref{ass:cont}}, for
every $\gamma>0$ the limit $\kappa_L(\delta,\gamma;C)$ exists and
\[
\kappa_L(\delta,\gamma;C)
\;=\;
\kappa_L(\delta,\tilde\gamma;C)
\;=\;
\tilde\kappa(\delta).
\]
In particular, $\gamma\mapsto\kappa_L(\delta,\gamma;C)$ is constant on
$(0,\infty)$.
\end{theorem}

After proving the theorem, we verify that the Pareto-mixture copula
$C_{\PM}$ defined in the problem statement satisfies the assumptions
of the theorem and hence inherits scaling invariance.

\section{Proof of Theorem~\ref{thm:main}}

Fix $\delta\in(0,2)$ and $\gamma>0$. We aim to show that, for every
$\eps>0$, there exists $u^\star=u^\star(\eps)\in(0,1)$ such that
\begin{equation}\label{eq:goal}
\bigl|E_\gamma(u;\delta)-E_{\tilde\gamma}(u;\delta)\bigr|
\;\le\;\eps
\qquad \text{for all }u\in(0,u^\star].
\end{equation}
Combined with $\lim_{u\downarrow 0}E_{\tilde\gamma}(u;\delta)=\tilde\kappa(\delta)$ from
Assumption~\ref{ass:limit}, this implies that $E_\gamma(u;\delta)$ also
converges to $\tilde\kappa(\delta)$ as $u\downarrow 0$, which is the
claim of the theorem.

\subsection*{Step 1: A key identity linking $\gamma$ to $\tilde\gamma$
via a $u$-dependent shift of $\delta$}

Set
\[
\Delta \;:=\; \log\!\bigl(\gamma/\tilde\gamma\bigr) \in \R.
\]
For each $u\in(0,1)$ write $t:=\log u<0$ and define
\begin{equation}\label{eq:du}
\delta_u \;:=\; \delta + \frac{\Delta}{t}
\;=\; \delta + \frac{\log(\gamma/\tilde\gamma)}{\log u}.
\end{equation}
Because $t\to-\infty$ as $u\downarrow 0$ while $\Delta$ is fixed, we
have
\begin{equation}\label{eq:du-conv}
\delta_u\;\longrightarrow\;\delta \qquad (u\downarrow 0).
\end{equation}
The point of \eqref{eq:du} is the following identity. Using $u=e^t$
and $\gamma/\tilde\gamma=e^\Delta$,
\[
\tilde\gamma\,u^{\delta_u}
\;=\;\tilde\gamma\,e^{(\delta+\Delta/t)t}
\;=\;\tilde\gamma\,e^{\delta t}\,e^{\Delta}
\;=\;\gamma\,u^{\delta},
\]
and similarly
\[
\tilde\gamma^{-1}u^{\,2-\delta_u}
\;=\;\tilde\gamma^{-1}e^{(2-\delta-\Delta/t)t}
\;=\;\tilde\gamma^{-1}e^{(2-\delta)t}\,e^{-\Delta}
\;=\;\gamma^{-1}u^{\,2-\delta}.
\]
Substituting into definition \eqref{eq:Edef}, the numerator and
denominator both match, so
\begin{equation}\label{eq:identity}
E_{\gamma}(u;\delta) \;=\; E_{\tilde\gamma}(u;\delta_u)
\qquad \text{for every }u\in(0,1).
\end{equation}
Identity \eqref{eq:identity} is the cornerstone of the argument:
evaluating the ratio at the scaling $\gamma$ along direction $\delta$
is exactly evaluating it at the reference scaling $\tilde\gamma$ along
the perturbed direction $\delta_u$, and the perturbation
$\delta_u-\delta=\Delta/\log u$ vanishes as $u\downarrow 0$.

\subsection*{Step 2: Triangle inequality and three pieces to control}

Fix $\eps>0$. From \eqref{eq:identity},
\[
E_\gamma(u;\delta) - E_{\tilde\gamma}(u;\delta)
\;=\;
E_{\tilde\gamma}(u;\delta_u) - E_{\tilde\gamma}(u;\delta).
\]
Insert $\pm\tilde\kappa(\delta_u)\pm\tilde\kappa(\delta)$ inside the
absolute value and apply the triangle inequality:
\begin{equation}\label{eq:triangle}
\bigl|E_{\tilde\gamma}(u;\delta_u)-E_{\tilde\gamma}(u;\delta)\bigr|
\;\le\;
\underbrace{\bigl|E_{\tilde\gamma}(u;\delta_u)-\tilde\kappa(\delta_u)\bigr|}_{(\mathrm{I})}
+
\underbrace{\bigl|\tilde\kappa(\delta_u)-\tilde\kappa(\delta)\bigr|}_{(\mathrm{II})}
+
\underbrace{\bigl|\tilde\kappa(\delta)-E_{\tilde\gamma}(u;\delta)\bigr|}_{(\mathrm{III})}.
\end{equation}
We show that each of the three terms is at most $\eps/3$ for all
sufficiently small $u$, each step justified by a different
assumption.

\subsection*{Step 3: Bounding each term by $\eps/3$}

\paragraph{Term (III): pointwise convergence at $\delta$.}
By Assumption~\ref{ass:limit}, $E_{\tilde\gamma}(u;\delta)\to\tilde\kappa(\delta)$ as
$u\downarrow 0$. Hence there exists $u_1\in(0,1)$ such that
\begin{equation}\label{eq:bound3}
\bigl|E_{\tilde\gamma}(u;\delta)-\tilde\kappa(\delta)\bigr|\;\le\;\eps/3
\qquad\text{for all }u\in(0,u_1].
\end{equation}

\paragraph{Term (II): continuity of $\tilde\kappa$ at $\delta$.}
By Assumption~\ref{ass:cont}, $\tilde\kappa$ is continuous at $\delta$,
so for the given $\eps>0$ there exists $\eta>0$ such that
$|\tilde\kappa(d)-\tilde\kappa(\delta)|\le\eps/3$ whenever
$d\in(0,2)$ and $|d-\delta|\le\eta$. By \eqref{eq:du-conv} there
exists $u_2\in(0,1)$ such that $|\delta_u-\delta|\le\eta$ for all
$u\in(0,u_2]$ (in particular $\delta_u\in(0,2)$ for $u$ small).
Therefore
\begin{equation}\label{eq:bound2}
\bigl|\tilde\kappa(\delta_u)-\tilde\kappa(\delta)\bigr|\;\le\;\eps/3
\qquad\text{for all }u\in(0,u_2].
\end{equation}

\paragraph{Term (I): local uniform convergence near $\delta$.}
This is the term that requires more than mere pointwise convergence:
the argument $\delta_u$ at which $E_{\tilde\gamma}(\cdot;\,\cdot)$ is
evaluated depends on $u$, so pointwise convergence at the fixed point
$\delta$ alone is insufficient. By
Assumption~\ref{ass:uniform} applied at $\delta_0=\delta$ with tolerance
$\eps/3$, there exist $r>0$ and $u_0\in(0,1)$ such that
\begin{equation}\label{eq:local-unif}
\bigl|E_{\tilde\gamma}(u;d)-\tilde\kappa(d)\bigr|\;\le\;\eps/3
\qquad
\text{for all }u\in(0,u_0]\text{ and all }d\in(0,2)\text{ with }
|d-\delta|\le r.
\end{equation}
Using \eqref{eq:du-conv} once more, there exists $u_3\in(0,1)$ such
that $|\delta_u-\delta|\le r$ for all $u\in(0,u_3]$. Combining with
\eqref{eq:local-unif} yields
\begin{equation}\label{eq:bound1}
\bigl|E_{\tilde\gamma}(u;\delta_u)-\tilde\kappa(\delta_u)\bigr|\;\le\;\eps/3
\qquad\text{for all }u\in(0,\min(u_0,u_3)].
\end{equation}

\subsection*{Step 4: Combining the bounds}

Let
\[
u^\star \;:=\; \min(u_0,u_1,u_2,u_3) \;\in\;(0,1).
\]
For every $u\in(0,u^\star]$, the three estimates
\eqref{eq:bound1}, \eqref{eq:bound2}, \eqref{eq:bound3} apply
simultaneously. Substituting into \eqref{eq:triangle},
\[
\bigl|E_{\tilde\gamma}(u;\delta_u)-E_{\tilde\gamma}(u;\delta)\bigr|
\;\le\;\frac{\eps}{3}+\frac{\eps}{3}+\frac{\eps}{3}\;=\;\eps.
\]
By identity \eqref{eq:identity}, this is precisely
\[
\bigl|E_{\gamma}(u;\delta)-E_{\tilde\gamma}(u;\delta)\bigr|\;\le\;\eps
\qquad\text{for all }u\in(0,u^\star],
\]
which is the claim \eqref{eq:goal}.

\subsection*{Step 5: Conclusion}

Since $\eps>0$ was arbitrary, \eqref{eq:goal} together with
$E_{\tilde\gamma}(u;\delta)\to\tilde\kappa(\delta)$ implies that
\[
\lim_{u\downarrow 0}E_\gamma(u;\delta)
\;=\;\lim_{u\downarrow 0}E_{\tilde\gamma}(u;\delta)
\;=\;\tilde\kappa(\delta).
\]
Hence $\kappa_L(\delta,\gamma;C)$ exists and equals
$\kappa_L(\delta,\tilde\gamma;C)$. As $\gamma>0$ and $\delta\in(0,2)$
were arbitrary, the angular lower-tail coefficient is independent of
the scaling parameter. \qed

\begin{remark}[The case $\tilde\kappa(\delta)=+\infty$]
If $\tilde\kappa(\delta)=+\infty$, then by definition, for every
$M>0$ there exists $u_M\in(0,1)$ such that
$E_{\tilde\gamma}(u;\delta)\ge M$ for all $u\in(0,u_M]$. The same
strategy applies with an ``$M$-large'' one-sided version of
Assumption~\ref{ass:uniform}: for every $M>0$ there exist $r>0$ and
$u_0\in(0,1)$ such that $E_{\tilde\gamma}(u;d)\ge M$ for all
$u\in(0,u_0]$ and all $d$ with $|d-\delta|\le r$. Together with
identity~\eqref{eq:identity} and $\delta_u\to\delta$, this gives
$E_\gamma(u;\delta)\ge M$ for all sufficiently small $u$, so
$\kappa_L(\delta,\gamma;C)=+\infty$ as well. The $\eps$-inequality in
Step~2 is simply replaced by this divergence statement.
\end{remark}

\section{Verification for the Pareto-Mixture Copula $C_{\PM}$}

We now check that the Pareto-mixture copula $C_{\PM}$ introduced in the
problem statement satisfies
Assumptions~\ref{ass:limit}--\ref{ass:cont} for every $\delta\in(0,2)$,
so that Theorem~\ref{thm:main} applies and yields scaling invariance.

\subsection{Recall of the model and marginal tails}

Recall that $Z\sim\Pa(\theta)$ means $\bar F_Z(z)=z^{-\theta}$ for
$z\ge 1$. Set
\[
X_1=\frac{R}{R^*}\cdot\frac{H_1}{H_1^*},\qquad
X_2=\frac{R}{R^*}\cdot\frac{H_2}{H_2^*},
\]
with $R\sim\Pa(\beta)$, $R^*\sim\Pa(\alpha)$, $H_i\sim\Pa(b_i)$,
$H_i^*\sim\Pa(a_i)$, all mutually independent and with strictly
positive shape parameters.

A standard Mellin-transform calculation shows that, for a ratio of two
independent Pareto factors, $\Prob(R/R^*\le x)\sim K_{\alpha,\beta}\,x^{\alpha}$ and
$\Prob(H_i/H_i^*\le y)\sim K_{a_i,b_i}\,y^{a_i}$ as the arguments tend to
$0$. By independence of $R/R^*$ and $H_i/H_i^*$, the lower tail of
$X_i$ is regularly varying:
\begin{equation}\label{eq:marginal}
\Prob(X_i\le x)\;\sim\; L_i\,x^{m_i}, \qquad x\downarrow 0,
\end{equation}
with $m_i=\min(\alpha,a_i)>0$ and $L_i>0$. Hence the quantile satisfies
\begin{equation}\label{eq:quantile}
F_{X_i}^{-1}(u)\;\sim\;\ell_i\,u^{1/m_i}, \qquad u\downarrow 0,
\end{equation}
for some $\ell_i>0$.

\subsection{Joint lower tail and verification of the assumptions}

By conditioning on $R/R^*$ and using the independence of $H_1/H_1^*$
and $H_2/H_2^*$, one obtains the joint lower-tail asymptotics in
homogeneous form: there exists a positively homogeneous function
$\rho:(0,\infty)^2\to(0,\infty)$ such that
\begin{equation}\label{eq:joint-rv}
\Prob(X_1\le x_1,\,X_2\le x_2)
\;=\;\rho(x_1,x_2)\,(1+o(1)),
\qquad (x_1,x_2)\to(0,0),
\end{equation}
where the $o(1)$ remainder is uniform on every compact set
of rays $\{(x_1,x_2):x_1=c_1 s,\;x_2=c_2 s\}$ with $(c_1,c_2)$ in a
compact subset of $(0,\infty)^2$.

Substituting \eqref{eq:quantile} into the definition of $C_{\PM}$ and
applying \eqref{eq:joint-rv} at
$x_i=F_{X_i}^{-1}\!\bigl(\gamma^{(-1)^{i+1}}u^{\delta_i}\bigr)$, with
$\delta_1=d$ and $\delta_2=2-d$, gives
\begin{equation}\label{eq:CPM-asymp}
C_{\PM}\!\bigl(\gamma u^{d},\,\gamma^{-1}u^{2-d}\bigr)
\;=\;G(\gamma,d)\,u^{\kappa^\star(d)}\bigl(1+o(1)\bigr),
\qquad u\downarrow 0,
\end{equation}
where $G(\gamma,d)>0$ is continuous in $d$ on $(0,2)$ and bounded away
from $0$ and $\infty$ on each compact subset of $(0,2)$, and where
$\kappa^\star:(0,2)\to(0,\infty)$ is a continuous function of the
direction $d$ alone, depending on the shape parameters but not on
$\gamma$. Moreover, the $o(1)$ remainder in \eqref{eq:CPM-asymp} is
uniform in $d$ on each compact subset of $(0,2)$, since the
remainder in \eqref{eq:joint-rv} is uniform on compact ray families and
the quantile asymptotics \eqref{eq:quantile} are uniform on compact
ranges of $\gamma^{\pm 1}u^{\delta_i}$.

Taking logarithms in \eqref{eq:CPM-asymp},
\begin{equation}\label{eq:CPM-log}
\log C_{\PM}\!\bigl(\gamma u^{d},\,\gamma^{-1}u^{2-d}\bigr)
\;=\;\kappa^\star(d)\log u + \log G(\gamma,d) + o(1),
\end{equation}
where the $o(1)$ is uniform in $d$ on compact subsets of $(0,2)$. Dividing by
$\log u$, which is negative for $u\in(0,1)$ and tends to $-\infty$,
gives
\begin{equation}\label{eq:E-CPM}
E_\gamma(u;d)
\;=\;\kappa^\star(d) + \frac{\log G(\gamma,d)}{\log u} + \frac{o(1)}{\log u}
\;=\;\kappa^\star(d) + o(1)
\qquad (u\downarrow 0),
\end{equation}
again uniformly in $d$ on each compact subset of $(0,2)$.

From \eqref{eq:E-CPM} we read off the three assumptions for
$C=C_{\PM}$ and any reference scale $\tilde\gamma>0$:
\begin{itemize}[noitemsep,topsep=2pt]
\item (Assumption~\ref{ass:limit}) For every $d\in(0,2)$ the pointwise
limit exists and equals $\tilde\kappa(d)=\kappa^\star(d)$, which is
finite.
\item (Assumption~\ref{ass:uniform}) The convergence in \eqref{eq:E-CPM}
is uniform on compact subsets of $(0,2)$; in particular, for any
$\delta\in(0,2)$, choosing a small closed interval $[\delta-r,\delta+r]\subset(0,2)$
gives the local uniform convergence required.
\item (Assumption~\ref{ass:cont}) $\tilde\kappa(d)=\kappa^\star(d)$ is
continuous on $(0,2)$, in particular at $\delta$.
\end{itemize}

\subsection{Conclusion for $C_{\PM}$}

\begin{corollary}\label{cor:pm}
For every $\delta\in(0,2)$ the angular lower-tail coefficient
$\kappa_L(\delta,\gamma;C_{\PM})$ of the Pareto-mixture copula is
well-defined as a finite limit and is independent of the scaling
parameter $\gamma>0$. Explicitly, $\kappa_L(\delta,\gamma;C_{\PM})=\kappa^\star(\delta)$,
which depends only on $\delta$ and the shape parameters
$(\alpha,\beta,a_1,a_2,b_1,b_2)$ of the underlying Pareto factors.
\end{corollary}

\begin{proof}
Combine the verification of
Assumptions~\ref{ass:limit}--\ref{ass:cont} above with
Theorem~\ref{thm:main}. Alternatively, \eqref{eq:E-CPM} already
yields the limit directly: as $u\downarrow 0$, the term
$\log G(\gamma,d)/\log u\to 0$ since the numerator is a fixed finite
number for each fixed $\gamma$ and $d$, and the denominator diverges to
$-\infty$. Hence $\lim_{u\downarrow 0}E_\gamma(u;\delta)=\kappa^\star(\delta)$ for
every $\gamma>0$.
\end{proof}

\section{Summary}

We have proved:
\begin{enumerate}[label=(\roman*),leftmargin=*]
\item Under the regularity conditions stated as
Assumptions~\ref{ass:limit}--\ref{ass:cont}---existence of the
$\kappa_L$-limit at a single reference scale $\tilde\gamma$, locally
uniform convergence in the direction parameter, and continuity of the
limit function at $\delta$---the angular lower-tail coefficient
$\kappa_L(\delta,\gamma;C)$ does not depend on $\gamma$
(Theorem~\ref{thm:main}). The proof rests on the explicit identity
$E_\gamma(u;\delta)=E_{\tilde\gamma}(u;\delta_u)$, where
$\delta_u=\delta+\log(\gamma/\tilde\gamma)/\log u\to\delta$, combined
with a triangle-inequality decomposition into pointwise convergence,
continuity, and local uniform convergence terms, each controlled to
$\eps/3$.
\item The Pareto-mixture copula $C_{\PM}$ satisfies
Assumptions~\ref{ass:limit}--\ref{ass:cont} for every $\delta\in(0,2)$
because its lower tail has a regular-variation expansion of the form
$C_{\PM}(\gamma u^{d},\gamma^{-1}u^{2-d})=G(\gamma,d)u^{\kappa^\star(d)}(1+o(1))$
with $\kappa^\star$ continuous, $G$ positive and finite, and remainder
uniform in $d$ on compact subsets of $(0,2)$. Hence
$\kappa_L(\delta,\gamma;C_{\PM})=\kappa^\star(\delta)$ does not depend
on $\gamma$ (Corollary~\ref{cor:pm}).
\end{enumerate}
Thus the angular lower-tail coefficient $\kappa_L(\delta,\gamma;C)$ is
genuinely a function of the direction $\delta$ alone (given $C$), as
the intuition motivating its definition required.

\end{document}